\section{Addition on \texorpdfstring{$\mathbb{N}$}{N}}

\begin{topicbox}{Recursive Definition of Addition}
\begin{exposition}
Addition on a Peano system is defined from the successor primitive by
recursion. For each fixed left input $m$, the map $n \mapsto m+n$ is generated
from the initial value $S(m)$ and the successor operator on the same Peano
system. This definition makes addition a derived operation rather than an
independent axiom.
\end{exposition}

\begin{tcolorbox}[colback=defbox, colframe=defborder, arc=2pt,
  left=6pt, right=6pt, top=4pt, bottom=4pt,
  title={\small\textbf{Definition (Addition on a Peano System)}},
  fonttitle=\small\bfseries]
\begin{definition}[Addition on a Peano System]
\label{def:addition-on-peano-system}
Let $(P,S,1)$ be a Peano system. For each $m \in P$, define the function
\[
f_m : P \to P
\]
to be the unique function satisfying
\[
f_m(1)=S(m)
\qquad\text{and}\qquad
\forall n \in P\;(f_m(S(n))=S(f_m(n))).
\]
For $m,n \in P$, define
\[
m+n \coloneqq f_m(n).
\]
\end{definition}
\end{tcolorbox}

\begin{remark*}[Standard quantified statement]
\[
\begin{aligned}
&\text{Let } (P,S,1) \text{ be a Peano system. For each } m \in P,
\text{ let } f_m : P \to P \text{ satisfy}\\
&f_m(1)=S(m)
\qquad\text{and}\qquad
\forall n \in P\;(f_m(S(n))=S(f_m(n))).\\
&\text{Then for all } m,n \in P,\quad m+n \coloneqq f_m(n).
\end{aligned}
\]
\end{remark*}

\begin{remark*}[Interpretation]
Addition is introduced by fixing the left input and recursively generating the
right-input map. The value $m+1$ is set to be the successor of $m$, and each
further successor step in the right input advances the output by one more
successor. This makes addition repeated succession.
\end{remark*}

\begin{remark*}[Dependencies]
\begin{itemize}
  \item \hyperref[def:peano-system]{Peano System}
  \item \hyperref[thm:peano-iterator-theorem]{Peano Iterator Theorem}
\end{itemize}
\end{remark*}

\begin{tcolorbox}[colback=thmbox, colframe=thmborder, arc=2pt,
  left=6pt, right=6pt, top=4pt, bottom=4pt,
  title={\small\textbf{Theorem (Addition Is Well-Defined on a Peano System)}},
  fonttitle=\small\bfseries]
\begin{theorem}[Addition Is Well-Defined on a Peano System]
\label{thm:addition-well-defined-on-peano-system}
Let $(P,S,1)$ be a Peano system. For every $m \in P$, there exists a unique
function
\[
f_m : P \to P
\]
such that
\[
f_m(1)=S(m)
\qquad\text{and}\qquad
\forall n \in P\;(f_m(S(n))=S(f_m(n))).
\]
Consequently the rule
\[
(m,n) \longmapsto m+n \coloneqq f_m(n)
\]
defines a unique binary operation on $P$.
\hyperref[prf-addition-well-defined-on-peano-system]{Go to proof.}
\end{theorem}
\end{tcolorbox}

\begin{remark*}[Standard quantified statement]
\[
\begin{aligned}
&\text{Let } (P,S,1) \text{ be a Peano system.}\\
&\forall m \in P\;\exists! f_m : P \to P\;
\Bigl(f_m(1)=S(m) \land
\forall n \in P\;(f_m(S(n))=S(f_m(n)))\Bigr).
\end{aligned}
\]
\end{remark*}

\begin{remark*}[Interpretation]
The theorem shows that the recursive clauses for addition determine exactly one
right-input map for each fixed left input. The addition operation is therefore
not assumed; it is constructed from the successor map by recursion.
\end{remark*}

\begin{remark*}[Dependencies]
\begin{itemize}
  \item \hyperref[def:addition-on-peano-system]{Addition on a Peano System}
  \item \hyperref[thm:peano-iterator-theorem]{Peano Iterator Theorem}
\end{itemize}
\end{remark*}
\end{topicbox}
